\begin{abstract}
本論文は、
自己と環境を区別し、
入力を受けて行動を選択し、
内部状態を遷移させ、
内部モデル状態を更新しながら存続を図る系を、
共通の形式で記述する対象として自律認知システム
（ACS：Autonomous Cognitive System）を定式化する。
本理論は著者が研究する倫理的協調知性理論
（ECIT：Ethical Coordinative Intelligence Theory）の基盤として、
倫理的協調を可行性保存の観点から形式的に検討するための
前段階の形式枠組みを与えることを目的とする。
本論文の形式的内容は
Phase~1（決定論的・離散時間で、状態・行動・内部モデル状態空間にコンパクト性を仮定する枠組み）
に限定され、
単一ACSの内部動学を中心に扱う。
組織ACSについては、
単一ACS枠組みを社会的組織へ構造的に適用する
限定的な記述として扱う。
ACSを7タプル $\Sigma=(X,U,E,f,\Pi,M,\mathrm{Update})$ として定義し、
持続評価関数 $\Pi$ から派生する
可行集合 $K:=\{x\in X\mid \Pi(x)\geq 0\}$ を可行性記述の中心的対象とする。
長期的な可行性の維持に関わる制約軸として、
自己維持（SM）・環境利用（EU）・可行応答（FR）の三原理を導入し、
それぞれの時刻別充足述語と論理的地位を区別して整理する。
FRは、必要な局面で可行な応答が得られること、
およびその必要応答ポテンシャルの保持・獲得・回復・再構成に関わる制約軸である。
これら三原理を原理水準の制約軸とし、
これとは別の機構水準の記述として、
可行性に関わる内部評価と応答の形成を担う
三機構情報処理アーキテクチャ
（TIPA：Triadic Information-Processing Architecture）、
その作動を支える物理的基盤層（PSL）および
実装部位層（ISL）からなる内部機構構造を与える。
あわせて、一段の作動記録が次の時刻へ閉じる標準作動形式と時刻間の接続、
および外部対象との関係性の最小構造と関係品質を定式化する。
さらに、可行行動が存在しない閉塞と、
構造的個体としての同一性を成立させていた組織関係が
同一個体として回復しない形で失われた崩壊とを区別する。
また、ACSが参加する集団・組織と
組織自身のACS性を区別し、
組織自身について整形式なACS構造が実現する場合を、
個体ACSの参加の有無によらず組織ACSとして記述し、先行節の定義を組織水準へ適用する。
ACS理論は、異なる自律主体の行動を、特定の文化的善悪判断に先立って、
可行性上の要請、内部状態、外部環境、関係および内部機構の作動との
構造的な関係として記述するための共通の形式を目指すものであり、
本論文はその決定論的・構造的な基礎を与える。
本論文は記述理論かつ構造理論の立場をとり、
人間、生物、AI、ロボット、組織等への
実証的妥当性を本論文の形式的結果から直接には主張しない。
決定論的枠組みのもとでの補助的詳細化はPhase~1.$\alpha$へ、
確率的拡張、連続時間化、因果的関係動学、
分岐を許す時間的自己同一性、
および社会動態の定量化はPhase~2へ、
ECITとの正式な接続はPhase~3へ委譲する。
\end{abstract}
\bigskip
\noindent\textbf{キーワード}：自律認知システム、可行性、Viability Theory、持続評価関数、三原理（自己維持・環境利用・可行応答）、可行性保存制約、TIPA、組織ACS、閉塞、崩壊。
